\clearpage
\fancyhead[L]{\small MERCURIUS IN SOLE VISUS.}
\hypertarget{native-653-lower}{}
\begin{center}\small Göttingische gelehrte Anzeigen. Stück 169. Seite 1688. 21. Oktober 1833.\\[3pt]\rule{.72\linewidth}{.3pt}\end{center}

\emph{Mercurius in Sole visus, of Overgang van Mercurius over de Zon; den 5 Mei \textup{1832} te Utrecht waargenomen door} G. \textsc{Moll}. Amsterdam 1833. Bei C. G. Sulpke. 43 Seiten. 4.

Wir finden in dieser kleinen Schrift die Beobachtungen des letzten Merkurdurchgangs in Utrecht auf der Sternwarte, und in der Wohnung des Herrn \textsc{Van Beek}, wie auch in Leiden, Amsterdam und Nimwegen; ferner eine Zusammenstellung aller Beobachtungen früherer Durchgänge, die in Holland oder von holländischen Astronomen gemacht worden sind; schließlich die in Utrecht und Leyden angestellten Beobachtungen der Bedeckung des Saturn durch den Mond am 8. Mai 1832. Eine ausführlichere Besprechung wäre hier überflüssig, da der wesentliche Inhalt aus Briefen des Herrn Professors \textsc{Moll} an Herrn Staatsrat \textsc{Schumacher} bereits im 229. Heft der Astronomischen Nachrichten bekannt gemacht worden ist.
